\documentclass[11pt]{article}
% =========================================================
%  學生講義 共用前言（以 XeLaTeX 編譯）
%  需要套件：xeCJK, tcolorbox（其餘多在 BasicTeX 內）
% =========================================================
\usepackage[a4paper,margin=2.3cm]{geometry}
\usepackage{amsmath,amssymb}
\usepackage{xcolor}
\usepackage{graphicx}

% ---- 中文字型（macOS 系統字型）----
\usepackage{xeCJK}
\setCJKmainfont{Songti TC}      % 內文：宋體
\setCJKsansfont{Heiti TC}       % 標題/強調：黑體
\xeCJKsetup{AutoFakeBold=2.2, AutoFakeSlant=0.2}

% ---- 版面 ----
\setlength{\parindent}{0pt}
\setlength{\parskip}{0.55em}
\linespread{1.15}
\renewcommand{\baselinestretch}{1.15}

% ---- 顏色 ----
\definecolor{cBlue}{HTML}{1f6feb}
\definecolor{cGray}{HTML}{6b7280}
\definecolor{cGrayBg}{HTML}{f3f4f6}
\definecolor{cPurp}{HTML}{7a3ff2}
\definecolor{cPurpBg}{HTML}{f5f1ff}

% ---- 標註框 ----
\usepackage[most]{tcolorbox}

% 就地補充新概念（灰底）：\begin{補充框}{標題}
\newtcolorbox{補充框}[1]{
  enhanced, breakable, arc=2pt, boxrule=0.6pt,
  colback=cGrayBg, colframe=cGray,
  left=9pt,right=9pt,top=7pt,bottom=7pt,
  before upper={\textbf{\sffamily #1}\par\vspace{3pt}},
}

% 延伸 / 開放問題（紫色左邊條）：\begin{延伸框}{標題}
\newtcolorbox{延伸框}[1]{
  enhanced, breakable, arc=2pt, boxrule=0pt,
  colback=cPurpBg, colframe=cPurpBg,
  borderline west={3pt}{0pt}{cPurp},
  left=10pt,right=9pt,top=7pt,bottom=7pt,
  before upper={\textbf{\sffamily 延伸閱讀｜#1}\par\vspace{3pt}},
}

% ---- 圖：置中、底下小字說明 ----
% \fig{檔名}{寬度}{說明}
\newcommand{\fig}[3]{%
  \begin{center}
  \includegraphics[width=#2\linewidth]{#1}\\[2pt]
  {\small\color{cGray} #3}
  \end{center}}

% ---- 章節標題用黑體 ----
\newcommand{\h}[1]{\sffamily #1}


\begin{document}

\begin{center}
{\sffamily\Huge\bfseries 數A4-3 機率統計}\\[3pt]
{\sffamily\large 普通高中 數學（數學A・第四冊）・學生講義}
\end{center}
\vspace{0.6em}

在第一冊的排列組合與機率中，學到的是「\textbf{算一個事件的機率}」：樣本空間、古典機率與期望值。本章往前一步，探討事件\textbf{彼此之間}的關係。

主線分三段。先把「機率到底是什麼意思」釐清，區分\textbf{主觀機率（subjective probability）}與\textbf{客觀機率（objective probability）}。接著進入\textbf{條件機率（conditional probability）}：「已知一件事發生後，另一件事的機率如何改變」，並由它導出\textbf{乘法公式（multiplication rule）}、事件的\textbf{獨立（independence）}，以及必須與「互斥」分清楚的觀念。最後把這些組合起來，得到\textbf{全機率公式（law of total probability）}與\textbf{貝氏定理（Bayes' theorem）}，學會把判斷從「先驗（prior）」修正到「後驗（posterior）」，並用它解釋為什麼「篩檢驗出陽性，不等於真的得病」。

條件機率與貝氏定理是現代資料科學、機器學習與醫學決策的共同語言，也是大學機率論與統計推論的起點。

\medskip
本章主線：

\begin{center}
兩種機率（主觀／客觀）$\rightarrow$ 條件機率 $\rightarrow$ 乘法公式 $\rightarrow$ 獨立性（與互斥區分）$\rightarrow$ 全機率公式 $\rightarrow$ 貝氏定理（篩檢偽陽性）
\end{center}

\section*{\sffamily 4-3-1\quad 主觀機率與客觀機率}

\subsection*{\sffamily A.\ 「機率是 0.7」到底在說什麼}

「機率」這個詞，其實有兩種\textbf{意思不同}的用法。

\begin{itemize}\setlength{\itemsep}{1pt}
\item 「明天\textbf{降雨機率} 70\%」——這是對一件\textbf{只會發生一次}的事（明天）所下的判斷。
\item 「這顆公正硬幣出現正面的機率是 $\frac12$」——這是因為硬幣兩面\textbf{對稱、等可能}，或因為「丟很多次，大約一半是正面」。
\end{itemize}

這兩種用法背後，是兩種不同的機率觀。

\begin{補充框}{兩種機率觀：客觀機率與主觀機率}
\textbf{客觀機率（objective probability）}：機率是事物本身的客觀性質，與「誰在猜」無關。它有兩個來源：
\begin{itemize}\setlength{\itemsep}{1pt}
\item \textbf{古典（對稱）觀}：當所有基本結果\textbf{等可能}時，
$$P(A)=\frac{\text{有利結果數}}{\text{全部結果數}}.$$
這就是第一冊學的古典機率。
\item \textbf{頻率（次數）觀}：在\textbf{相同條件下重複大量試驗}，事件 $A$ 出現的\textbf{相對次數}會逐漸穩定，這個穩定值即為 $P(A)$：
$$P(A)\approx \frac{A\ \text{出現的次數}}{\text{試驗總次數}}\quad(\text{試驗次數很大時}).$$
\end{itemize}

\textbf{主觀機率（subjective probability）}：機率是「\textbf{某個人對某件事的相信程度}」，用 0 到 1 的數字表示（0 表示完全不信會發生，1 表示確信會發生，0.5 表示一半一半）。「明天降雨 70\%」「某球隊奪冠機率三成」都屬此類——這些事\textbf{無法重複試驗}，但人仍須對它下判斷。
\end{補充框}

\subsection*{\sffamily B.\ 由數據獲得客觀機率（相對次數）}

當一件事沒有明顯的對稱性可用（例如一顆做工不均的圖釘，丟下去尖朝上的機率），便無法用古典機率直接計算。這時只能靠\textbf{做實驗、數次數}：把圖釘丟很多次，記錄尖朝上的累計相對次數。隨著試驗次數增加，相對次數會\textbf{逐漸穩定}在某個值附近，這個穩定值即為機率的估計。試驗次數越多，估計越可靠。

\textbf{注意：}試驗次數太少時，相對次數會劇烈跳動，不能當作機率。例如「丟 10 次有 7 次正面，所以正面機率 0.7」並不成立——十次的結果只是一次抽樣，運氣成分很大。頻率觀的機率是\textbf{試驗次數趨於很大}時相對次數的穩定值（大數法則）。

\fig{d43-relative-frequency}{0.86}{圖 4-3-1：相對次數的穩定。模擬擲一顆「尖朝上機率 0.63」的圖釘，橫軸為試驗次數、縱軸為累計相對次數。次數少時（左側陰影區）相對次數劇烈跳動、不能當作機率；隨著次數增加，相對次數逐漸穩定在 0.63（紅虛線），這個穩定值才是頻率觀下的機率。}

\subsection*{\sffamily C.\ 用機率的性質檢驗主觀機率}

主觀機率雖然是「個人的相信程度」，但\textbf{不能任意給定}：只要它要被當成機率使用，就\textbf{必須遵守機率的基本性質}，否則會自相矛盾。

\begin{補充框}{機率公理：任何機率都要滿足}
設樣本空間為 $S$，事件 $A,B\subseteq S$：
$$0\le P(A)\le 1,\qquad P(S)=1,\qquad P(\varnothing)=0.$$
兩事件互斥（$A\cap B=\varnothing$，不會同時發生）時：$P(A\cup B)=P(A)+P(B)$。由此可得\textbf{餘事件}與\textbf{加法定理}：
$$\boxed{\,P(A^c)=1-P(A)\,}\qquad \boxed{\,P(A\cup B)=P(A)+P(B)-P(A\cap B)\,}.$$
\end{補充框}

舉例而言，若某球評對甲、乙、丙三隊奪冠給的主觀機率分別是 $0.5$、$0.4$、$0.3$，由於奪冠是互斥事件且三隊之外無人能奪冠，三者機率\textbf{總和應為 1}；但 $0.5+0.4+0.3=1.2>1$，違反「總機率為 1」，因此這組判斷\textbf{不合理}，至少要調整其中一個數字。

客觀機率與主觀機率並非對立，而是\textbf{回答不同問題的工具}：可重複、有對稱性的事（骰子、抽球）用客觀機率；一次性、無法重複的事（明天、選舉、某病人）只能用主觀機率，但它仍須服從機率公理。本章 4-3-3 的貝氏統計，正是把主觀機率（先驗）拿來，用數據（客觀）一步步修正。

\begin{延伸框}{主觀機率「憑感覺講的數字」，憑什麼也叫機率？}
客觀機率有「等可能」或「重複試驗」撐腰，主觀機率兩者都沒有。關鍵在於：\textbf{一個數字能否被當成「機率」，重點不是它怎麼來的，而是它有沒有遵守機率的規則。} 主觀機率雖然來源是個人判斷，但只要服從機率公理（每個機率在 $0$ 到 $1$ 之間、互斥又窮盡的可能機率總和為 1、包含關係一致），就能和客觀機率一樣參與運算（加法、乘法、貝氏更新）。

主觀機率為何非得服從公理不可，有一個很實際的理由：\textbf{如果主觀機率不一致，別人可以設計一組賭局，讓你無論結果如何都穩輸}。這在機率論裡稱為「荷蘭賭（Dutch book）」。換句話說，服從機率公理不只是數學要求，也是不被套利的理性底線。這也說明合理的主觀機率，與客觀機率服從\textbf{完全相同}的規則。

常見誤解是「主觀就是不準、客觀就是準」。兩者回答的是\textbf{不同問題}，沒有誰比較準的問題。氣象的「降雨 70\%」其實是用大量歷史相似天氣資料校準過的，相當可靠；要評估它準不準，須看「所有報 70\% 的日子裡，是否約七成真的下雨」。
\end{延伸框}

\section*{\sffamily 4-3-2\quad 條件機率與獨立性}

\subsection*{\sffamily A.\ 已知一個情報，機率就改變}

擲一顆公正骰子，「出現 6 點」的機率是 $\frac16$。但若\textbf{已知「這次擲出的是偶數」}，「是 6 點」的機率還是 $\frac16$ 嗎？

不是了。既然已知是偶數，可能的結果只剩 $\{2,4,6\}$ 三個，而 6 是其中之一，機率變成 $\frac13$。

關鍵轉變在於：一旦「已知 $B$ 發生」，考慮的範圍就\textbf{從整個樣本空間 $S$ 縮小成 $B$}。原本要在 $S$ 裡算比例，現在改在 $B$ 裡算比例。這個「在 $B$ 之內、$A$ 佔多少」的機率，就是\textbf{條件機率}。

\begin{補充框}{定義：條件機率}
在「事件 $B$ 已發生」的條件下，事件 $A$ 發生的機率，記作 $P(A\mid B)$（讀作「在 $B$ 之下 $A$ 的機率」）。當 $P(B)>0$ 時，
$$\boxed{\;P(A\mid B)=\frac{P(A\cap B)}{P(B)}\;}.$$
直覺：分母從整個 $S$ 換成 $B$（新的「全部」），分子是「同時在 $A$ 又在 $B$」的部分（$A\cap B$，新世界裡屬於 $A$ 的部分）。
\end{補充框}

\fig{d43-conditional}{0.92}{圖 4-3-2：條件機率把分母從 $S$ 換成 $B$。左——原本在整個樣本空間裡看 $A$；右——已知 $B$ 發生後世界縮小成 $B$，只剩 $A\cap B$ 這塊算數，$P(A\mid B)$ 即 $A\cap B$ 在 $B$ 裡所佔的比例。}

\begin{補充框}{為什麼分母換成 $P(B)$、分子換成 $P(A\cap B)$}
把樣本空間想成一整片草地，$A$、$B$ 是其上兩塊（可能重疊）的區域。

\textbf{還沒有任何情報時}，「$A$ 的機率」就是「$A$ 這塊地佔整片草地多少」。\textbf{現在已知「$B$ 發生」}，表示真正的結果一定落在 $B$ 裡，$B$ 以外的部分全部可以劃掉，世界縮小成只剩 $B$。在這個小世界裡再問「$A$ 的機率」，就是問「在 $B$ 之內、屬於 $A$ 的部分佔多少」，而「在 $B$ 裡又屬於 $A$」的部分正是重疊 $A\cap B$。因此
$$P(A\mid B)=\frac{A\cap B\ \text{的面積}}{B\ \text{的面積}}=\frac{P(A\cap B)}{P(B)}.$$
\begin{itemize}\setlength{\itemsep}{1pt}
\item \textbf{分母換成 $P(B)$}：新世界的「全部」就是 $B$，要拿它當分母，使新世界的總機率重新變回 1（$P(B\mid B)=\frac{P(B)}{P(B)}=1$）。
\item \textbf{分子換成 $P(A\cap B)$}：新世界裡能算數的，只有「同時也在 $B$ 裡」的那部分 $A$，落在 $B$ 外的 $A$ 已被劃掉。
\end{itemize}
除法的本質就是「重新標準化」：把舊世界的機率，按新世界的大小 $P(B)$ 重新縮放回「總和為 1」。
\end{補充框}

\textbf{注意：}條件機率不一定比原機率小。情報可能讓機率變大、變小或不變——例如已知「是偶數」後，「是 6 點」從 $\frac16$ 升到 $\frac13$（變大）。另外要區分 $P(A\mid B)$ 與 $P(A\cap B)$：後者是「在原本的大世界裡，$A$、$B$ 同時發生」的機率；前者是「已假定 $B$ 發生後」$A$ 的機率，分母不同，數值通常不同。

更重要的是，\textbf{條件不能隨意反過來}：一般情況下
$$P(A\mid B)\neq P(B\mid A).$$
例如「是 6 點」在「偶數」裡佔 $\frac13$，但「偶數」在「6 點」裡佔 1（已知是 6 點，當然是偶數），兩者差很多。這個方向性，正是 4-3-3 貝氏定理要處理的核心。

\subsection*{\sffamily B.\ 乘法公式}

把條件機率的定義移項，就得到一個算「\textbf{兩件事接連都發生}」的工具。

\begin{補充框}{定義：機率的乘法公式}
由 $P(A\mid B)=\dfrac{P(A\cap B)}{P(B)}$ 移項：
$$\boxed{\;P(A\cap B)=P(B)\,P(A\mid B)=P(A)\,P(B\mid A)\;}.$$
文字版：「$A$、$B$ 都發生」的機率 ＝「先發生其中一件」乘上「在那之後另一件也發生」。推廣到三件事（連續抽取最常用）：
$$P(A\cap B\cap C)=P(A)\,P(B\mid A)\,P(C\mid A\cap B).$$
這條由兩次乘法公式串成：把 $A\cap B$ 當成一個整體，先對「$(A\cap B)$ 與 $C$」用一次、再對「$A$ 與 $B$」用一次，即 $P(A\cap B\cap C)=P(A\cap B)\,P(C\mid A\cap B)=P(A)\,P(B\mid A)\,P(C\mid A\cap B)$，相當於「沿一條路徑把每一步的條件機率連乘」，抽 $n$ 次可一路乘下去。
\end{補充框}

乘法公式最典型的用途是\textbf{不放回連抽}。例如袋中 5 紅 3 白共 8 顆，不放回連抽兩顆，「兩顆都紅」的機率為
$$P(\text{第一紅})\times P(\text{第二紅}\mid\text{第一紅})=\frac58\times\frac47=\frac{5}{14}.$$
重點在於：不放回時，第二次的機率\textbf{依賴}第一次的結果（已抽走一紅，剩 4 紅、共 7 顆），所以第二次一定要用\textbf{條件機率}$\frac47$，不能寫成 $\frac58$。

\fig{d43-multiplication}{0.86}{圖 4-3-3：乘法公式的樹狀圖（袋中 5 紅 3 白，不放回連抽兩顆）。沿每條路徑把分岔機率連乘得葉子機率；通往「兩顆都紅」的路徑為 $\frac58\times\frac47=\frac{5}{14}$。第二步分母由 8 變成 7，正說明不放回時前後\textbf{不獨立}。}

\subsection*{\sffamily C.\ 事件的獨立性}

上面「不放回」時，第一次的結果\textbf{改變}了第二次的機率（$\frac58\to\frac47$）。但如果是\textbf{放回}再抽，第二次抽之前袋子完全恢復原狀，「第一次抽到什麼」對第二次\textbf{毫無影響}。這種「一件事發生與否，不改變另一件事的機率」，就是\textbf{獨立}。

\begin{補充框}{定義：兩事件獨立}
若「$B$ 是否發生」\textbf{不影響} $A$ 的機率，即 $P(A\mid B)=P(A)$，則稱 $A$、$B$ \textbf{獨立（independent）}。把它代回乘法公式 $P(A\cap B)=P(B)P(A\mid B)$，得到\textbf{對稱、好記、也是正式定義}的判準：
$$\boxed{\;A,B\ \text{獨立}\iff P(A\cap B)=P(A)\,P(B)\;}.$$
直觀：獨立時「兩件都發生」的機率，就是各自機率\textbf{直接相乘}（不必管誰先誰後、誰影響誰）。
\end{補充框}

由此可看出「放回 vs 不放回」的差別：同樣 5 紅 3 白連抽兩顆求「兩顆都紅」，\textbf{放回}時兩次獨立，$P=\frac58\times\frac58=\frac{25}{64}$；\textbf{不放回}時兩次不獨立，$P=\frac58\times\frac47=\frac{5}{14}$。兩者不同。判斷一題該怎麼乘，先問自己：\textbf{第二次的機率有沒有被第一次改變？}

\subsection*{\sffamily D.\ 獨立不等於互斥（最容易混淆，務必分清）}

「獨立」和「互斥」是最常被混為一談的兩個詞，但它們講的是\textbf{完全不同}的事，甚至幾乎\textbf{相反}：

\begin{itemize}\setlength{\itemsep}{1pt}
\item \textbf{互斥（mutually exclusive）}：兩件事\textbf{不會同時發生}，$A\cap B=\varnothing$，即 $P(A\cap B)=0$。問的是「能不能\textbf{同時}發生」。
\item \textbf{獨立（independent）}：一件事發生\textbf{不影響}另一件事的機率，$P(A\cap B)=P(A)P(B)$。問的是「會不會\textbf{互相影響}」。
\end{itemize}

\fig{d43-independence}{0.92}{圖 4-3-4：互斥與獨立的對照。左——互斥的兩圓\textbf{完全分開}（$A\cap B=\varnothing$，$P(A\cap B)=0$）；右——獨立的兩圓\textbf{相交}，重疊比例恰好等於各自比例相乘（$P(A\cap B)=P(A)P(B)\neq0$）。}

\begin{補充框}{為什麼互斥的兩個正機率事件「反而不獨立」}
關鍵在於：\textbf{互斥其實是一種最強烈的「互相影響」}。若 $A$、$B$ 互斥，那麼「$B$ 一旦發生，$A$ 就絕不可能發生」。用條件機率寫出來：
$$P(A\mid B)=\frac{P(A\cap B)}{P(B)}=\frac{0}{P(B)}=0.$$
本來 $A$ 有 $P(A)>0$ 的機率，一聽到「$B$ 發生」就\textbf{直接掉到 0}——這影響大得不能再大，怎麼可能是「互不影響」（獨立）？

用數學確認：兩個\textbf{正機率}事件 $A$、$B$ 若互斥，則
$$P(A\cap B)=0,\qquad P(A)P(B)>0,$$
兩者不相等，所以\textbf{不獨立}。一句話記住：\textbf{「互斥」是不會一起來；「獨立」是來不來互不干涉。兩個正機率事件若互斥，就一定不獨立。}

\textbf{唯一的退化例外：}若其中一個事件機率為 0（如不可能事件），則 $P(A\cap B)=0=P(A)P(B)$ 同時成立，此時互斥與獨立可兼得；但這只是「機率為 0」的瑣碎情形，理解主旨即可。
\end{補充框}

要小心兩個常見誤解：其一是「獨立 ＝ 沒有交集」——也錯，獨立的兩圓\textbf{一定有交集}（$P(A\cap B)=P(A)P(B)>0$），沒有交集（互斥）的反而不獨立。其二是把互斥當成獨立——剛好相反，兩個正機率事件互斥就保證不獨立。實務上判斷一題屬於哪一種，可循下列線索：

\begin{itemize}\setlength{\itemsep}{1pt}
\item 「\textbf{抽完不放回}、抽走就少一個」$\Rightarrow$ 前後\textbf{不獨立}（會互相影響）。
\item 「\textbf{放回}、兩顆不同的骰子或硬幣、兩個無關的人」$\Rightarrow$ 多半\textbf{獨立}。
\item 「\textbf{同一次試驗的不同結果}（如出現 1 或出現 2）」$\Rightarrow$ \textbf{互斥}。
\item 不確定時回到定義硬算：互斥看 $A\cap B$ 是不是空；獨立看 $P(A\cap B)$ 是否等於 $P(A)P(B)$。
\end{itemize}

實際上，大多數真實事件\textbf{既不互斥也不獨立}（兩圓有交集，但重疊比例不等於各自比例相乘），這是最常見的情形。

\section*{\sffamily 4-3-3\quad 貝氏定理}

\subsection*{\sffamily A.\ 全機率公式：先分類，再加總}

很多問題裡，一個結果可以\textbf{來自好幾種不同的「來源」}，而每個來源出這個結果的機率不一樣。例如一間工廠的產品來自機台 A 或 B，兩台的不良率不同，問「隨手抽一件，是不良品的機率」。

做法是把整個樣本空間\textbf{依來源切成幾塊}（彼此互斥、合起來涵蓋全部，稱為一組\textbf{分割，partition}），分別算每塊貢獻的機率，再加起來。

\begin{補充框}{定義：全機率公式}
設 $B_1,B_2,\dots,B_n$ 把樣本空間\textbf{互斥且窮盡}地分割（任兩不相交、合起來是全部），每個 $P(B_i)>0$。則對任一事件 $A$：
$$\boxed{\;P(A)=\sum_{i=1}^{n}P(B_i)\,P(A\mid B_i)\;}.$$
兩塊的版本（最常用）：
$$P(A)=P(B)\,P(A\mid B)+P(B^c)\,P(A\mid B^c).$$
直覺：$A$ 發生的總機率 ＝ 把「經由每一條來源路徑到達 $A$ 的機率」全部加起來（樹狀圖中每條路徑相乘、各路徑相加）。
\end{補充框}

\fig{d43-bayes-tree}{0.78}{圖 4-3-5：全機率公式的樹狀圖。先依來源分岔（機台 A 佔 0.6、機台 B 佔 0.4），每條路徑把分岔機率相乘，再把所有通往「不良」的路徑相加，得到總不良率 $0.6\times0.05+0.4\times0.10=0.07$。}

以圖 4-3-5 的工廠為例：60\% 來自機台 A、40\% 來自機台 B，A 的不良率 5\%、B 的不良率 10\%。設 $D$＝「不良」，則
$$P(D)=P(A)P(D\mid A)+P(B)P(D\mid B)=0.6\times0.05+0.4\times0.10=0.07.$$
整體不良率 7\%，介於 5\% 與 10\% 之間，且偏向佔比較大的 A，合理。

\subsection*{\sffamily B.\ 貝氏定理：把條件反過來}

上面算的是「\textbf{已知來自某機台}，是不良品的機率」（$P(D\mid A)$，由因到果）。但實務上常拿到的是\textbf{相反的問題}：「\textbf{已知抽到的是不良品}，它來自機台 A 的機率是多少」（$P(A\mid D)$，由果到因）。

前面已說明 $P(A\mid B)\neq P(B\mid A)$，所以\textbf{不能}把 $P(D\mid A)$ 直接當成 $P(A\mid D)$。把兩個方向接起來的橋樑，就是\textbf{貝氏定理}。

\begin{補充框}{定義：貝氏定理}
由乘法公式 $P(A\cap D)=P(A)P(D\mid A)=P(D)P(A\mid D)$，移項即得：
$$\boxed{\;P(A\mid D)=\frac{P(A)\,P(D\mid A)}{P(D)}\;}.$$
把分母 $P(D)$ 用全機率公式展開（兩塊版）：
$$P(A\mid D)=\frac{P(A)\,P(D\mid A)}{P(A)\,P(D\mid A)+P(A^c)\,P(D\mid A^c)}.$$
\begin{itemize}\setlength{\itemsep}{1pt}
\item $P(A)$：\textbf{先驗機率（prior）}——還沒看到證據 $D$ 之前，對 $A$ 的判斷。
\item $P(D\mid A)$：\textbf{似然（likelihood）}——「假設 $A$ 成立、會看到證據 $D$」的機率，是反方向（因$\rightarrow$果）的條件機率，通常由題目給定（如機台不良率、檢驗敏感度）。
\item $P(A\mid D)$：\textbf{後驗機率（posterior）}——看到證據 $D$ \textbf{之後}，修正過的判斷。
\end{itemize}
貝氏定理就是「看到證據後，如何把先驗更新成後驗」的公式，可記成口訣：\textbf{後驗 $\propto$ 先驗 $\times$ 似然}（分母 $P(D)$ 只是把總機率標準化回 1）。
\end{補充框}

承上工廠例：已知抽到的是不良品，求它來自機台 A 的機率。先驗 $P(A)=0.6$，似然 $P(D\mid A)=0.05$，分母 $P(D)=0.07$，故
$$P(A\mid D)=\frac{0.6\times0.05}{0.07}=\frac{0.03}{0.07}=\frac37\approx0.43.$$
機台 A 雖佔產量 6 成（先驗 0.6），但因不良率低，「\textbf{不良品裡}來自 A 的」只佔約 43\%——看到「不良」這個證據後，對 A 的判斷從 0.6 降到 0.43。

\subsection*{\sffamily C.\ 經典應用：疾病篩檢的偽陽性}

貝氏定理最有名、也最違反直覺的應用，是\textbf{醫學篩檢}：一個準確率聽起來很高的檢驗，驗出陽性時，受檢者\textbf{真正得病的機率卻可能很低}。

考慮一種疾病，人群中的\textbf{盛行率為 1\%}。某檢驗的\textbf{敏感度（sensitivity）99\%}（得病者驗出陽性的機率）、\textbf{特異度（specificity）95\%}（沒病者驗出陰性的機率）。設 $S$＝「得病」、$T$＝「驗出陽性」，已知
$$P(S)=0.01,\quad P(S^c)=0.99,\quad P(T\mid S)=0.99,\quad P(T\mid S^c)=1-0.95=0.05.$$
先用全機率公式算「驗出陽性的總機率」：
$$P(T)=P(S)P(T\mid S)+P(S^c)P(T\mid S^c)=0.01\times0.99+0.99\times0.05=0.0594.$$
再用貝氏定理算後驗：
$$P(S\mid T)=\frac{P(S)P(T\mid S)}{P(T)}=\frac{0.0099}{0.0594}\approx0.167.$$
即使驗出陽性，\textbf{真正得病的機率也只有約 16.7\%}；換句話說，陽性者裡超過八成是虛驚一場（偽陽性）。

\begin{補充框}{用「自然頻率」想最清楚：為什麼答案這麼低}
把抽象的機率換成\textbf{具體人數}就一目了然。假設 10000 人受檢：
\begin{itemize}\setlength{\itemsep}{1pt}
\item 真正得病的有 $10000\times1\%=100$ 人，其中 $100\times0.99\approx99$ 人驗出陽性（\textbf{真陽性}）。
\item 沒病的有 9900 人，其中 $9900\times5\%=495$ 人\textbf{也}驗出陽性（\textbf{偽陽性}）。
\item 驗出陽性的共 $99+495=594$ 人，真得病的只有 99 人：$\dfrac{99}{594}\approx16.7\%$。
\end{itemize}
\textbf{關鍵}：健康的人\textbf{基數太大}（9900 人），即使偽陽性率只有 5\%，乘上 9900 也產生 495 個偽陽性，遠多於 99 個真陽性。盛行率（先驗）越低，這個「基數效應」越嚴重，後驗就被壓得越低。
\end{補充框}

\fig{d43-screening}{0.62}{圖 4-3-6：10000 人受檢的自然頻率方塊圖（每格 1 人）。真陽性 99 人、偽陽性 495 人、真陰性 9405 人、偽陰性 1 人。驗出陽性者共 594 人，真得病者只有 99 人，故 $P(\text{病}\mid\text{陽})=99/594\approx17\%$。}

\textbf{注意：}不要把「敏感度」當成「得病機率」。$99\%$ 是 $P(T\mid S)$（得病者中陽性的比例，是檢驗本身的性質）；受檢者真正在意的是 $P(S\mid T)$（陽性者中真得病的比例）。兩者一個是「病 $\rightarrow$ 陽」、一個是「陽 $\rightarrow$ 病」，\textbf{方向相反}，數值可以差很多。把這兩個搞混，稱為\textbf{基本比率謬誤（base rate fallacy）}——忽略了盛行率這個基本比率，是連受過訓練的醫師都常犯的錯誤。

\begin{延伸框}{貝氏定理在哲學上仍有兩派（但不影響高中計算）}
貝氏定理本身是\textbf{數學恆等式}，由乘法公式移項而得，必然成立，沒有任何爭議。有不同學派的是它背後的機率哲學：

\begin{itemize}\setlength{\itemsep}{1pt}
\item \textbf{頻率學派（frequentist）}：機率只能談「可重複試驗的相對次數」，反對把「一個人是否得病」這種一次性命題賦予機率，因此對「主觀先驗」抱持保留。
\item \textbf{貝氏學派（Bayesian）}：機率就是「相信程度」，可以對任何命題給先驗，再用貝氏定理隨證據更新。現代機器學習、醫療決策、垃圾郵件過濾大多採貝氏思路。
\end{itemize}

這是統計學的方法論之爭，\textbf{至今並存}，各有適用場景。但無論站哪一派，上面篩檢例子的算術結果完全相同。此外要說明：篩檢得到約 16.7\% 並\textbf{不代表這種檢驗沒用}——它把得病機率從 1\% 提升到 16.7\%（約十幾倍），足以篩出該做進一步確診檢查的對象。篩檢的定位本就是「初篩」而非「確診」。
\end{延伸框}

\subsection*{\sffamily D.\ 先驗到後驗：機率是會更新的判斷}

回頭看 4-3-1 的主觀機率。貝氏定理給了主觀機率一個\textbf{動態的生命}：先憑經驗給一個\textbf{先驗}（相信程度），每當有\textbf{新證據}進來，就用貝氏定理把它更新成\textbf{後驗}；後驗又可以當成下一輪的先驗，繼續被新證據修正。這就是現代\textbf{貝氏統計（Bayesian statistics）}的核心精神，也是垃圾郵件過濾、醫療診斷、機器學習裡無所不在的思路。

由此把全章串成一條主線：古典機率告訴你還沒看資料前該信多少（先驗），條件機率提供「已知一個情報後如何更新」的工具，貝氏定理則告訴你看到資料後該改成多少（後驗）。\textbf{機率不是冰冷的數字，而是會隨證據更新的判斷。}

\section*{\sffamily 4-3-4\quad 本章重點整理}

\begin{補充框}{一頁複習（公式與關鍵結論）}
\begin{itemize}\setlength{\itemsep}{2pt}
\item \textbf{兩種機率}：客觀（古典＝對稱等可能、頻率＝大量試驗的相對次數穩定值）／主觀（個人相信程度，仍須服從機率公理）。
\item \textbf{機率公理}：$0\le P\le1$、$P(S)=1$、互斥可加；$P(A^c)=1-P(A)$；$P(A\cup B)=P(A)+P(B)-P(A\cap B)$。
\item \textbf{條件機率}：$P(A\mid B)=\dfrac{P(A\cap B)}{P(B)}$（樣本空間縮小成 $B$）；一般 $P(A\mid B)\neq P(B\mid A)$。
\item \textbf{乘法公式}：$P(A\cap B)=P(B)P(A\mid B)=P(A)P(B\mid A)$（連抽、不放回必用）。
\item \textbf{獨立}：$P(A\cap B)=P(A)P(B)\iff P(A\mid B)=P(A)$（互不影響；放回 $\Rightarrow$ 獨立）。
\item \textbf{獨立 $\neq$ 互斥}：互斥是 $A\cap B=\varnothing$（不會同時發生）；\textbf{兩個正機率事件若互斥，必然不獨立}。
\item \textbf{全機率公式}：$P(A)=\sum_i P(B_i)P(A\mid B_i)$（先依來源分類，再加總）。
\item \textbf{貝氏定理}：$P(A\mid D)=\dfrac{P(A)P(D\mid A)}{P(D)}$；先驗 $\rightarrow$ 後驗。\textbf{篩檢陽性 $\neq$ 真得病}：盛行率低時偽陽性壓垮真陽性（基本比率謬誤）。
\end{itemize}
\end{補充框}

\end{document}
